% =============================================================================
% 6. CONCLUSIONES
% =============================================================================
\section{Conclusiones}
\label{sec:conclusions}

\subsection{Resumen de logros}

Se logr\'o desplegar exitosamente el modelo Chest Xpert AI como un servicio containerizado de producci\'on que cumple todos los criterios de la r\'ubrica:

\begin{enumerate}
    \item \textbf{Imagen Docker funcional}: arranca con un solo \texttt{docker run}, sin dependencias externas ni vol\'umenes.
    \item \textbf{Modelo funcional}: responde con predicciones correctas y pertinentes para las 5 patolog\'ias tor\'acicas.
    \item \textbf{Documentaci\'on clara}: README con instrucciones de construcci\'on, ejecuci\'on y ejemplo completo.
    \item \textbf{Pipeline CI/CD}: linting, testing, build y publicaci\'on autom\'atica en GHCR.
    \item \textbf{Versionado sem\'antico}: releases autom\'aticos basados en Conventional Commits.
    \item \textbf{Observabilidad}: stack LGTM completo con traces, logs, m\'etricas y profiling continuo.
\end{enumerate}

\subsection{Decisiones t\'ecnicas clave}

\begin{itemize}
    \item \textbf{ONNX Runtime sobre PyTorch}: reducci\'on de dependencias de $\sim$2GB a $\sim$50MB, manteniendo la misma precisi\'on del modelo con inferencia 10x m\'as r\'apida.
    \item \textbf{uv sobre pip}: resoluci\'on de dependencias 10--100x m\'as r\'apida, lockfile determinista que garantiza reproducibilidad bit-a-bit.
    \item \textbf{Modelo embebido en la imagen}: sacrifica tama\~no de imagen por simplicidad operativa. El profesor no necesita configurar nada.
    \item \textbf{Multi-stage build}: minimiza la superficie de ataque y el tama\~no final de la imagen al excluir herramientas de compilaci\'on.
\end{itemize}

\subsection{Lecciones aprendidas}

\begin{enumerate}
    \item \textbf{El shebang importa}: en builds multi-stage, los scripts del virtualenv quedan con rutas del builder stage. Soluci\'on: usar \texttt{python -m modulo} en vez de ejecutar el script directamente.
    \item \textbf{GITHUB\_TOKEN es suficiente para GHCR}: no se necesitan tokens personales para publicar im\'agenes en el Container Registry del mismo repositorio.
    \item \textbf{Semantic-release + GITHUB\_TOKEN no dispara workflows}: por dise\~no de GitHub. Los tags deben re-pushearse desde local o usar un PAT para triggear el CD.
    \item \textbf{Pinear versiones de Actions}: usar \texttt{v8} puede fallar si no existe el major tag; \texttt{v8.2.0} es m\'as seguro.
\end{enumerate}

\subsection{Trabajo futuro}

\begin{itemize}
    \item \textbf{Dashboards preconfigurados}: exportar dashboards de Grafana como JSON provisioning para despliegue autom\'atico.
    \item \textbf{Alertas}: configurar alertas en Grafana cuando la latencia p95 supere un umbral o el error rate suba.
    \item \textbf{Frontend containerizado}: la aplicaci\'on Angular servida desde Nginx como segunda imagen Docker.
    \item \textbf{Kubernetes}: despliegue con Deployment + Service + Ingress para escalabilidad horizontal.
    \item \textbf{Monitoreo de data drift}: alertas autom\'aticas cuando la distribuci\'on de entradas cambia respecto al training set.
    \item \textbf{Endpoint de explicabilidad}: Grad-CAM bajo demanda para interpretar predicciones individuales.
\end{itemize}
